\documentclass[10pt,a4paper]{amsart}
\usepackage[margin=19mm]{geometry}
\usepackage{amsmath,amssymb,mathtools}
\usepackage[scr=rsfs,cal=euler]{mathalfa}
\usepackage{xfrac}
\usepackage[unicode,hidelinks,bookmarks=false]{hyperref}
\newcommand{\ie}{i.e.\ }
\newcommand{\cf}{cf.\ }
\newcommand{\resp}{respectively\ }
\newcommand{\Subsection}{\subsection}
\providecommand{\nobreakdash}{\nobreak\hbox{-}}
\setlength{\emergencystretch}{2em}
\multlinegap0pt
\usepackage{bm}
\usepackage{bbm}
\usepackage{dsfont}
\usepackage{xspace}
\usepackage{cancel}
\usepackage{mathtools}
\usepackage{amsthm}
\makeatletter
\@ifundefined{theorem}{\newtheorem{theorem}{Theorem}}{}
\makeatother
\newcommand{\bI}{\mathbf{I}}%
\newcommand{\bW}{\mathbf{W}}%
\newcommand{\ba}{\boldsymbol{\alpha}}%
\newcommand{\bb}{\mathbf{b}}%
\newcommand{\be}{\mathbf{e}}%
\newcommand{\bun}{\mathbf{1}}%
\newcommand{\bz}{\mathbf{z}}%
\title[The Well-Tempered Classifier: Some Elementary Properties of ]{The Well-Tempered Classifier: Some Elementary Properties of Temperature Scaling}
\author{Pierre-Alexandre Mattei, Bruno Loureiro}
\date{Source: arXiv:2602.14862v2 (CC BY 4.0)}
\begin{document}
\maketitle

\noindent\emph{Source and licence.} The Well-Tempered Classifier: Some Elementary Properties of Temperature Scaling. Version arXiv:2602.14862v2 (CC BY 4.0), distributed under the Creative Commons Attribution 4.0 International licence, as recorded in the arXiv licence metadata for this exact version and shown on the abstract page https://arxiv.org/abs/2602.14862v2. Full rights notice: licenses/arxiv-2602.14862-CC-BY-4.0.txt.\par\smallskip

\noindent\textit{Editorial scope.} Counted unit: the Theorem whose source label is \texttt{eq:invarianceapp} in the article (source0.tex lines 569--579, source label \texttt{eq:invarianceapp}), delivered together with the article's own complete original proof of it (source0.tex lines 581--625, one \texttt{proof} environment of 45 lines). No mathematics is written, repaired or re-derived: statement and proof are the article's own text line for line. Required context retained uncounted: none. Every definition, display and label of those lines is retained unchanged. Omitted from this package and not redistributed: 1--568, 580--580, 626--647. Nothing inside a retained excerpt is omitted or altered: every retained body is delivered exactly as the article's lines are written, so no omission receipt of any kind accompanies this package. Citations: the retained excerpts carry no citation command at all, so no bibliography entry of the article is transcribed and no reference is redistributed. Reference adaptation: none was needed and none was made; every reference inside the delivered unit resolves inside it, so no unresolved reference or citation is produced. Printed slips kept literal: no printed word, symbol, hypothesis or number of the counted statement or of its proof was altered. Layout and class: the article's own class is \texttt{article} with packages including inputenc, fontenc, url, geometry, mathtools, amsmath, amssymb, amsthm, bbm, physics, colorblind, graphicx, hyperref, caption; the wrapper is the lane's pinned \texttt{amsart} editorial wrapper at 10pt on A4, which restates the theorem-like environments the retained text needs and adds the article's own macro definitions that the retained text needs (\texttt{\textbackslash bI}, \texttt{\textbackslash bW}, \texttt{\textbackslash ba}, \texttt{\textbackslash bb}, \texttt{\textbackslash be}, \texttt{\textbackslash bun}, \texttt{\textbackslash bz}), with extra wrapper packages (bm, bbm, dsfont, xspace, cancel, mathtools). The delivered numbering is the unit's own. Assets: no retained span contains a figure environment, a graphics-inclusion command, a PDF-page inclusion, a table or a TikZ picture; no image file is copied into this package, the package declares no asset dependency and no image file is redistributed with it. Licence: version arXiv:2602.14862v2 of this article is distributed under the Creative Commons Attribution 4.0 International licence, as recorded in the arXiv licence metadata for this exact version and shown on the abstract page https://arxiv.org/abs/2602.14862v2. A full rights notice is delivered at licenses/arxiv-2602.14862-CC-BY-4.0.txt. Characters: every retained excerpt is pure ASCII, so the delivered engine, XeLaTeX with its default Latin Modern text font, reports no missing character.
\begin{theorem}
    Let $\bW \in \mathbb{R}^{K \times K}$ and  $\bb \in \mathbb{R}^K$ such that, for all $\bz \in \mathbb{R}^K$ we have
    \begin{equation}
    \label{eq:invarianceapp}
        \textup{argmax}(\bW \bz+\bb) = \textup{argmax}(\bz).
    \end{equation}
Then, there exist $\ba \in \mathbb{R}^K$, $\beta> 0$, and $\gamma \in \mathbb{R}$ such that
\begin{equation}
    \bW = \beta \bI_K + \bun_K \ba ^T, \text{ and } \bb = \gamma \bun_K.
\end{equation}
\end{theorem}

\begin{proof}
The roadmap of the proof goes as follows: first, we will prove that $\bb$ has the desired form, then we will prove that $\bW$ has the desired form (by studying separately the cases $K=2$ and $K \geq 3$). Finally, we will show that $\beta>0$.



Applying Equation \eqref{eq:invarianceapp} to $\bz=0$ leads to $\text{argmax}(\bb)= \text{argmax}(\boldsymbol{0}_K)= \{1, \ldots, K \} $, therefore all coefficients of $\bb$ need to be equal, and thus there exists $\gamma \in \mathbb{R}$ such that $\bb = \gamma \bun_K$.

We now turn to $\bW$. The fact that $\bb = \gamma \bun_K$ implies that, for all $\bz \in \mathbb{R}^K$, 
\begin{equation}
\label{eq:invariance_linear}
    \text{argmax}(\bW \bz+\bb)= \text{argmax}(\bW \bz) = \text{argmax}(\bz),
\end{equation}
which will be our main tool to study the structure of $\bW$.

We begin with the case $K=2$. Let $\bW \in \mathbb{R}^{2 \times 2}$ that verifies \eqref{eq:invariance_linear}, and let $\alpha_1 = w_{21}$, $\alpha_2 = w_{12}$, and $\beta = w_{11}- \alpha_1$. The fact that $\text{argmax}(\bW \bun_2) = \text{argmax}(\bun_2)= \{1,2\}$ implies that
\begin{equation}
    w_{11} + w_{12} = w_{21} +w_{22} \implies w_{22} = w_{11} + w_{12} - w_{21} = \beta + \alpha_1 + \alpha_2 - \alpha_1 = \beta + \alpha_2.
\end{equation}
Thus, $\bW = \beta \bI_K + \bun_K \alpha ^T$. 




Let us now treat the case $K\geq 3$. For any $j \in \{1,\ldots,K\}$, let us denote $\be_j\in \mathbb{R}^K$ the vector whose only nonzero coefficient is the $j$th one, which is equal to $1$. Let $\bW$ be a matrix that satisfies \eqref{eq:invariance_linear}. Let $i,j,k$ be three distinct indices in $\{1, \ldots, K\}$. Applying the condition to the vector $\be_i + \be_j + \be_k$ leads to
\begin{equation}
    w_{ii} + w_{ij} + w_{ik} = w_{ji} + w_{jj} + w_{jk} = w_{ki} + w_{kj} + w_{kk},
\end{equation}
and applying it to  $\be_i + \be_k$ leads to
\begin{equation}
    w_{ii} + w_{ik} = w_{ki} + w_{kk}.
\end{equation}
Subtracting these equations implies that $w_{ij} = w_{kj}$. This means that, for each column $j$, all elements but the diagonal one will have the same value, that we will denote $\alpha_j$. 

Let us now use condition \eqref{eq:invariance_linear} with the vector $\bun_K$. This implies that 
\begin{equation}
    w_{11} + \sum_{j \neq 1} \alpha_j =  w_{22} + \sum_{j \neq 2} \alpha_j = \ldots = w_{KK} + \sum_{j \neq K} \alpha_j.
\end{equation}
Subtracting $\alpha_1 + \ldots + \alpha_K$ from all terms leads then yields
\begin{equation}
    w_{11} - \alpha_1 = \ldots =  w_{KK} - \alpha_K.
\end{equation}
Let us denote this quantity $\beta$. Putting the pieces together, we have proven that $\bW = \beta \bI_K + \bun_K \ba ^T$.

The only thing left to prove is the positivity of $\beta$. This follows from the fact that $\text{argmax}(\bW \be_1) = \text{argmax}(\be_1) = \{1\}$. Indeed, this implies that $w_{11}>w_{21}$ i.e. that $\beta + \alpha_1 > \alpha_1$, leading to $\beta>0$.
\end{proof}

\end{document}
